% Folder 93, batch 12 — archivists' pages 221–235.
% Pass: Opus 5.5 (claude-opus-5-5), 2026-10-10 — first pass, unchecked against the pages by a human.
%
% All fifteen pages are manuscript notes in Grothendieck's hand; none is
% skipped. Three runs: pp. 221–229, integrals of closed relative 1-forms along
% paths joining two sections, and the condition for an expression
% sum D_i int omega_i to be a Picard-Fuchs equation (equations numbered (1) to
% (15), (13 bis)); pp. 230–232, differential operators with constant
% coefficients on a stratified module (items 1) to 3)); pp. 233–235, his own
% pages « 1 » and « 2 », construction of the canonical connection on De Rham
% cohomology by torsors, which runs on past the end of the batch.
% The script letter he uses on pp. 223–227 for f_*(Z Omega^1)/d f_*(O_X) is
% set as \mathscr{F}; the reading of the letter is noted once.
\documentclass[11pt,a4paper]{article}
\input{../preamble/grothendieck}

\folder{93}
\batch{12}
\pages{221}{235}
\foldertitle{Connexions de Gauss-Manin et équations de Picard-Fuchs. Opération de Cartier : copies de tapuscrit annoté (s.d.), tiré à part (1981), notes manuscrites (s.d.).}
\dating{[1972]-1981}
\watermark{Édition de démonstration}

\begin{document}

\section{Intégrales de formes fermées relatives et équations de Picard-Fuchs}

\page{221}
1) Soit $f : X \to S$ un morphisme lisse d'espaces analytiques \add{[lisses]}
(holomorphe), $\omega$ une 1-forme fermée \add{relative sur $X/S$}, $a$ et $b$
deux sections \add{holomorphes} de $X$ sur $S$, telles que $\forall s \in S$,
$a(s)$ et $b(s)$ soient dans la même comp.\ connexe de $X_s$. Localement sur
$S$, on peut trouver un arc \struck{diff} \add{(différentiable par morceaux)}
$C_s : [0,1] = I \to X_s$, joignant $a(s)$ et $b(s)$, et dépendant
continûment de $s$.

Utilisant le fait que $\omega$ est fermée, donc localement de la forme
$d_{X/S}\varphi$, on constate que la fonction sur $S$
\[
  I(s) = \int_{C_s} \omega_s = \int_{a(s)}^{b(s)} \omega_s
\]
est holomorphe en $s$. [\add{plus précisément au voisinage de $s_0 \in S$ et
\ill{}} \uncertain{En effet}, \uncertain{supposons} \ill{} \ill{} \struck{$a(s_0)$}
$a(s_0)$ et $b(s_0)$ sont assez voisins pour que dans un ouvert $U$ contenant
$a(s_0)$, $b(s_0)$, on ait $\omega|U = d\varphi$, donc pour $s$ voisin de
$s_0$, $I(s) = \varphi(b(s)) - \varphi(a(s))$].
\uncertain{Changeant} le \ill{} des $C_s$ variant : le remplaçant par
$C'_s = C_s + Z_s$, $Z_s$ un 1-cycle, donc on trouve
\[
  I'(s) = I(s) + \int_{Z_s} \omega_s . \tag{1}
\]

Considérons une expression formelle (avec des $\omega_i$ comme $\omega$, et des
op.\ diff.\ \add{(rel.)} $D_i$ sur $S$]
\[
  \sum D_{i,s} \int_{a(s)}^{b(s)} \omega_{i,s} . \tag{2}
\]

\page{222}
Cette fonction est déterminée, indépendamment du choix de $C_s$, modulo une
expression de la forme
\[
  \sum D_{i,s} \int_{Z_s} \omega_{i,s}
  \quad \text{\struck{$= \int_{Z_s} \sum D_{i,s}\, \omega_{i,s}$}} \tag{3}
\]
\struck{[N.B. Le fait que $Z_s$ est un 1-cycle permet}
Supposons que les $D_i$, $\omega_i$ soient tels que l'expression précédente
soit nulle, quelle que soit la famille continue $s \mapsto Z_s$ de 1-cycles
\struck{$\Sigma$} \uncertain{définie} sur un ouvert de $S$ :
\struck{Alors l'expression} \add{On dira} que \struck{les}
$\mu = (D_i, \omega_i)_{i \in I}$ \emph{définit une équation de Picard-Fuchs}
sur $X/S$. Alors pour deux sections hol.\ quelc.\ $a$, $b$ de $X/S$,
l'expression
\[
  \mu(a,b) = \sum_i D_{i,s} \int_{a(s)}^{b(s)} \omega_{i,s} \tag{4}
\]
est bien déterminée, et est une fonction holomorphe sur $S$. Bien sûr on aura
\[
  \mu(a,b) + \mu(b,c) = \mu(a,c) . \tag{5}
\]
Cherchons une condition suffisante maniable, \add{sous forme différentielle,}
\struck{\add{\ill{}} \ill{} infinitésimale, et non plus globale,} pour que
$\mu$ soit de Pic.\ Fuchs. Notons

\page{223}
que sur
\[
  \mathscr{F} = f_*(Z\Omega^1_{X/S}) / d f_*(\mathcal{O}_X) , \tag{6}
\]
\note{la lettre cursive notée ici $\mathscr{F}$ (pp. 223 à 227) est de forme incertaine}
il y a une connexion intégrable naturelle, grâce à la formule de Cartan …
\marginal{attention, en réalité il faut que $X \to S$ soit \uncertain{steinien}}
Les \uncertain{coeff.} \uncertain{définissent} les sections $\omega'_i$ de
$\mathscr{F}$, et l'on peut former les $D_i \omega'_i$. Je dis qu'il suffit
qu'on ait
\[
  \sum D_i \omega'_i = 0 . \tag{7}
\]
Lorsque $X/S$ est muni d'une connexion intégrable, alors le calcul de
$\sum D_i \omega'_i$ \add{dans $f_*(\Omega^*_{X/S})$ sera aussi stratifié,
donc} se fait ainsi : \add{\uncertain{on dit que}} $\overline{D}_i$ se
\uncertain{relevant} \struck{pour la stratification en} \add{opérateurs}
$\overline{D}_i$ sur $\underline{\Omega}^*_X$ \add{commutant à la diff.\ rel.\
$d_{X/S}$} et \struck{\ill{}} \struck{qui à une forme fermée] \ill{} une forme}

\note{le cadre qui suit, numéroté (8) et (9), est barré de grands traits en croix}
\struck{(8)} $\alpha_{\overline{D}_i}\, \omega_i
= \alpha_{\overline{D}{}^0_i}\, \omega_i + c_i\, \omega_i$ ;
$\overline{D}_i = \overline{D}{}^0_i + c_i\, \mathrm{id}$
[$\overline{D}{}^0_i$ à coeff.\ cts, $c_i$ section de $\mathcal{O}_X$] ;
\struck{(9)} $\alpha_{\overline{D}_i}\, d\varphi
= \alpha_i\, \overline{D}_i \varphi$
\note{tout ce cadre est biffé ; la lettre lue $\alpha$ et le dernier membre de (9) sont incertains}

Considérons alors
\[
  \sum \overline{D}_i\, \omega_i = \varpi \tag{8}
\]
qui est encore une forme fermée. On vérifie alors \add{aisément} que l'on a

\page{224}
\[
  \sum D_i\, \omega'_i = \varpi' \quad \text{dans } \mathscr{F} \tag{9}
\]
et la condition (7) signifie que l'on a, localement au dessus de $S$, une
relation
\[
  \varpi \;(= \textstyle\sum \overline{D}_i\, \omega_i) = d\varphi ,
  \qquad \varphi \in \Gamma(X, \mathcal{O}_X) . \tag{10}
\]
Pour voir que cette relation implique bien la nullité de (3), il suffit de
prouver
\[
  \sum_i D_{i,s} \int_{Z_s} \omega_{i,s} = \int_{Z_s} \varpi \tag{11}
\]
qui implique, si $\varpi = d\varphi$, que le premier membre est égal à
$\int_{\partial Z_s} \varphi = 0$ car $\partial Z_s = 0$.

En fait, (11) provient de la relation générale
\[
  D_s \int_{Z_s} \omega_s = \int_{Z_s} \overline{D}_s\, \omega_s \tag{12}
\]
valable pour tout op.\ diff.\ hol.\ $D$ sur $S$, et toute forme
\add{holomorphe fermée} $\omega$ sur $X$ [Le fait que $\omega$ est fermée
implique que $\int_{Z_s} \omega$ est fonction holomorphe en $s$, et donne un
sens au premier membre de (12)]. Pour prouver cette

\page{225}
dernière, on [peut se ramener, en \uncertain{décomposant}
$Z_{s_0}$ en \add{petits} \struck{\ill{}} chemins $a^0_i \to a^0_{i+1}$, et
transportant \struck{par} horizontalement les $a^0_i$ \ill{}] à prouver la
relation plus générale
\note{une longue ligne ondulée traverse les deux premières lignes de la page ; biffure possible du passage entre crochets}
\[
  D_s \int_{a(s)}^{b(s)} \omega_s = \int_{a(s)}^{b(s)} (\overline{D}\omega)_s ,
  \tag{13}
\]
valable pour des chemins d'intégration $a(s)$, $b(s)$, qui sont des
\emph{sections horizontales} pour la connexion. [N.B. Bien entendu, on
suppose, pour donner un sens à cette formule, qu'on a choisi les arcs $C_s$
\uncertain{comme dessus}, donc (13) s'écrit avec plus de précision
\[
  D_s \int_{C_s} \omega_s = \int_{C_s} (\overline{D}\omega)_s
  \tag{13 bis}
\]
].

Pour prouver (13), on est ramené, en \add{se plaçant au voisinage de $s_0$ et}
décomposant le chemin $\int_{a(s)}^{b(s)}$, au cas où $a(s_0)$, $b(s_0)$ sont
assez voisins, de sorte qu'on peut se placer dans un ouvert $U$ où
\uncertain{l'on a} $\omega = d\varphi$. Mais alors le 1er
membre est $D_s\bigl(\varphi(b(s)) - \varphi(a(s))\bigr)$, le

\page{226}
deuxième est
\[
  \int_{a(s)}^{b(s)} \overline{D}\, d\varphi
  = \int_{a(s)}^{b(s)} d\, \overline{D}\varphi
  = (\overline{D}\varphi)(b(s)) - (\overline{D}\varphi)(a(s)) ,
\]
\struck{et les} et les deux expressions sont égales, grâce au fait que $a$,
$b$ sont des sections horizontales. Notons aussi que inversement, si $a$, $b$
sont des sections \emph{horizontales}, on aura
\[
  \mu(a,b) = \int_{a(s)}^{b(s)} \Bigl(\sum \overline{D}_i\, \omega_i\Bigr)_s
  = \varphi(b(s)) - \varphi(a(s)) , \tag{14}
\]
où $\varphi$ est défini par la relation (10).

Prouvons la suffisance de (7), pour la nullité de (3) en général. On va
prouver « plus généralement » qu'on aura alors
\[
  \sum D_{i,s} \int_{C_s} \omega_{i,s} = 0
\]
lorsque $C_s$ est une 1-chaîne dans $X_s$, variant différentiablement avec
$s$, dont le bord est « horizontal » en $s$, ou ce qui revient au même que
\[
  \sum_i D_{i,s} \int_{a(s)}^{b(s)} \omega_{i,s} = 0
\]
si $a$, $b$ sont des sections \emph{horizontales} de $X$, \add{sur un
\uncertain{ouvert} de $S$.} \struck{Pour} Se plaçant au voisinage d'un

\page{227}
point $s_0$ de $S$, on est ramené au cas où $a(s_0)$ est proche de $b(s_0)$,
de sorte qu'on peut \struck{\ill{}} se placer dans un ouvert de $X$ assez petit
pour qu'il y ait une connexion. \uncertain{Et on gagne}.

On peut aussi noter, \struck{quitte} pour prouver la \struck{formule} nullité
de (3) comme conséquence de (7), que l'on a
\[
  \sum D_{i,s} \int_{Z_s} \omega_i
  = \int_{Z_s} \Bigl(\sum \overline{D}_i\, \omega'_i\Bigr)_s
\]
\note{le numéro de cette formule est raturé}
où le deuxième membre est nul, car $\int_{Z_s} \varpi_s = 0$ si $\varpi$ est
un bord, et $(\sum \overline{D}_i\, \omega'_i)_s$ est une forme
\uncertain{définie modulo} un bord. La formule envisagée \struck{\ill{}}
\add{\ill{}}
\[
  \boxed{\; D_s \int_{Z_s} \omega'_s = \int_{Z_s} (\overline{D}\omega')_s \;}
  \tag{15}
\]
qui donne la propriété transcendante essentielle de la connexion intégrable
canonique sur $\mathscr{F}$. Nous n'allons pas prouver (15) ici, et
l'admettrons, mais notons que

\page{228}
(15) nous donne une \emph{réciproque} pour la condition envisagée pour que
$(D_i, \omega_i)$ soit de Picard-Fuchs : il faut et suffit que l'on ait (7).
Cela résulte du th.\ de De Rham sur les fibres $X_s$ ….

\emph{Calcul des $\mu(a,b)$} par voie « algébrique », sans calcul
d'intégrales. \add{C'est local sur $S$.} \add{La formule (5) nous} ramène à
calculer $\mu(a,b)$ \add{au voisinage de $s_0$,} pour \struck{$a(s)$}
$a(s_0)$ et $b(s_0)$ assez voisins. Cela nous permet alors de supposer que sur
$X/S$ il y a une connexion intégrable. On peut même supposer, par un
\uncertain{choix} convenable, que les sections $a$, $b$ sont horizontales.
Alors on trouve (14)
\[
  \mu(a,b) = \varphi(b(s)) - \varphi(a(s))
\]
où $\varphi$ est défini par (10).

Il y aurait lieu d'expliciter aussi une formule explicite, sans supposer que
$a(s)$ et $b(s)$ soient horizontales. Faisons le calcul au voisinage de
$s_0$, et soient $a_0$, $b_0$ les

\page{229}
sections horizontales au voisinage de $s_0$, telles que
\[
  a_0(s_0) = a(s_0) , \qquad b_0(s_0) = b(s_0) .
\]
On a
\[
  \mu(a,b) = \mu(a_0, b_0) + \mu(b_0, b) - \mu(a_0, a)
\]
\note{sous $\mu(a_0,b_0)$, un renvoi raturé, puis « (14) $\varphi(b_0) - \varphi(a_0)$ »}
donc
\begin{align*}
  \mu(a,b)(s_0) &= \bigl[\varphi(b(s_0)) - \varphi(a(s_0))\bigr] \\
  &\quad + \Bigl(\sum_i D_{i,s} \int_{b_0(s)}^{b(s)} \omega_{i,s}
  - \sum_i D_{i,s} \int_{a_0(s)}^{a(s)} \omega_{i,s}\Bigr)_{s = s_0}
\end{align*}
Il faut calculer les deux termes de \ill{} \ill{}, c'est une histoire
purement locale.

\note{la suite de la page est barrée de grands traits obliques}
\struck{Notons, grâce à la connexion de $X/S$,}
\[
  \text{\struck{$d_a \in \Gamma(S, a^*(\mathscr{V}_{X/S}) \otimes \Omega^1_S)$}}
\]
\struck{d'où, \struck{sur} contractant avec}
\[
  \text{\struck{$a^*(\omega) \in \Gamma(S, a^*(\Omega^1_{X/S}))$}}
\]
\struck{\add{dual de $a^*(\mathscr{V}_{X/S})$} une section, que nous notons
$a^*(\omega)$,}
\[
  \text{\struck{$a^*(\omega) \in \Gamma(S, \Omega^1_S)$}}
\]
\struck{Cette section} \struck{\add{forme} est encore \emph{fermée}, [N.B.
\uncertain{pour le cas où} \add{\uncertain{on l'étend}}, définissons
$a^*(\omega)$ pour toute forme ($\omega$ de degré quelconque) sur $X$, et
notons que, fermée ou non, $a^*$ commute à la différentielle ext.)}
\marginal{donner une formule …!}
\note{la lettre lue $\mathscr{V}_{X/S}$ (faisceau tangent relatif, semble-t-il) et le symbole lu $d_a$ sont incertains}

\section{Opérateurs différentiels à coefficients constants}

\page{230}
\marginal{$S \to T$ schéma relatif]}
1) Soit $F \xrightarrow{D} G \xrightarrow{D'} H$ une suite d'op.\ différentiels
\add{[d'ordre $1$]} qui est « formellement exacte ». Alors pour tout $P$, la
suite
\[
  \mathrm{Diff}^i(P,F) \xrightarrow{u \mapsto D \circ u}
  \mathrm{Diff}^{i+1}(P,G) \xrightarrow{v \mapsto D' \circ v}
  \mathrm{Diff}^{i+2}(P,H)
\]
est exacte, et la suite
\[
  \mathrm{Diff}^i(H,P) \xrightarrow{u \mapsto u \circ D'}
  \mathrm{Diff}^{i+1}(G,P) \xrightarrow{v \mapsto v \circ D}
  \mathrm{Diff}^{i+2}(F,P)
\]
est exacte.

2) Supposons que $E$ sur $S$ soit loc.\ libre \emph{stratifié}, et que le lemme
de Poincaré formel soit vrai pour le complexe
\[
  E \otimes \Omega^*_{S/T} = \Omega^*_{S/T}(E)
\]
[p.ex.\ $S$ lisse sur $T$, et caractéristiques de $T$ nulles]. Alors \add{$\forall
P$ quasi-coh.\ et $i \geq 0$} la suite
\begin{align*}
  0 \leftarrow \mathrm{Hom}(E,P) &\leftarrow \mathrm{Diff}^i(E,P)
  \xleftarrow{u \mapsto u \circ d^0_E} \\
  &\qquad \mathrm{Diff}^{i-1}(E \otimes \Omega^1_{S/T}, P)
  \xleftarrow{u \mapsto u \circ d^1_E} \cdots
\end{align*}
\note{au-dessus de la première flèche, une étiquette raturée suivie de « $(1)$ »}
[$\mathrm{Diff}^i(E,P) \simeq \mathrm{Hom}(P^i(E),P)
\xrightarrow{u \mapsto u \circ \delta_i} \mathrm{Hom}(E,P)$, où
$\delta_i : E \to P^i(E)$ défini par la stratification] est exacte.

\emph{N.B.} $\mathrm{Ker}\bigl(\mathrm{Diff}^i(E,P) \to \mathrm{Hom}(E,P)\bigr)$
est désigné par $\mathrm{Diff}^i_+(E,P)$, op.\ diff.\ « à coeff.\ cts » nul.\ ;
la stratification

\page{231}
Donc on trouve

\emph{Corollaire}
\[
  \mathrm{Diff}^{(i)}_+(E,P) \simeq
  \mathrm{Diff}^{i-1}(E \otimes \Omega^1_{S/T}, P) /
  \mathrm{Im}\, \mathrm{Diff}^{i-2}(E \otimes \Omega^2_{S/T}, P)
\]
Par suite, on a un homom.\ naturel injectif (?)
\[
  \mathrm{Diff}^{(i)}_+(E,P) \hookrightarrow
  \mathscr{H}\!om\,\mathrm{add}(\mathrm{Ker}\, d^1_E, P)
\]
\add{l'image est \uncertain{apparemment} : faisceau des op.\ diff.\ de
$\mathrm{Ker}\, d^1_E$ dans $P$}

Pour $i$ fixé, le $P$ « universel » pour les op.\ diff.\ à coeff.\ constants
(d'ordre $\leq i$) sur $E$ est
\[
  P^i(E) / \mathcal{O}_S \cdot \delta_i(E) ,
\]
où
\[
  \delta_i : E \to P^i(E)
\]
est défini par la stratification.

3) Supposons maintenant donné un sous-Module $N$ de $E$, on s'intéresse aux
homom.\ « différentiels » $\mathrm{Ker}\, d^1_E \to P$ qui s'annulent sur
$d^0_E N$, qui correspondent donc aux \struck{\ill{}} op.\ diff.\
$E \to P$, à coeff.\ cts, et qui s'annulent sur $N$, donc aux op.\ diff.\
$M = E/N \to P$ qui \struck{sont \ill{}} « proviennent » à coeff.\ cts sur
$E$. Ils correspondent \add{(pour l'ordre $\leq i$), aux homom.} linéaires
\[
  P^i(M) \to P
\]
qui s'annulent sur $\rho_i E$, où
\[
  \rho_i : E \to P^i(M)
\]

\page{232}
est l'homom.\ composé [linéaire pour la structure \emph{gauche} sur
$P^i(M)$]
\[
  E \xrightarrow{\delta_i} P^i(E) \to P^i(M) .
\]
Donc, pour $P$ variable, c'est représenté par
\[
  \mathscr{P}_i = P^i(M) / \mathcal{O}_S \cdot \rho_i(E) ,
\]
les $\mathscr{P}_i$ forment un syst.\ projectif strict, déterminé par la
connaissance de $M$ et de
\[
  \rho_\infty : E \to P^\infty(M)
\]
(linéaire pour la structure \emph{gauche} de $P^\infty(M)$), défini par les
$\rho_i$. Attention, ici $\rho_\infty$ satisfait à des conditions spéciales :
son composé $E \to M$ \struck{\ill{}} avec l'augmentation $P^\infty(M) \to M$
est un épimorphisme, + conditions d'intégrabilité ….

\section{Construction de la connexion canonique en cohomologie de De Rham}

\page{233}
\note{p. « 1 » de l'auteur ; le titre de section est le sien, encadré en tête de page}
\begin{tikzcd}
  X_0 \arrow[d, "f_0"] & X' \arrow[d, "f'"] & X'' \arrow[dl, "{f''}"] \\
  T_0 \arrow[r, hook] & T &
\end{tikzcd}

$T_0 \hookrightarrow T$ immersion 1-nilpotente, $T_0 = V(\underline{I})$ ;
\[
  X' \times_T T_0 \simeq X_0 , \qquad X'' \times_T T_0 \simeq X_0 ,
\]
$X'$, $X''$ deux prolongements lisses de $X_0/T_0$ sur $T$.
\marginal{schéma relatif formellement lisse (pour $f_0$) ; et, dans une bulle
reliée à $f_0$ : conditions de \uncertain{lissité} \add{de $X_0$, $X'$, $X''$}
satisfaites par $X'$, $X''$ \uncertain{qui sont} loc.\ $T$-isom.\ (sur $X_0$).}

$\mathscr{P}$ faisceau sur esp.\ $X_0$ des germes des $T$-isom.\
$X' \to X''$ ;

$\mathscr{G}$ — $T$-autom.\ $X'' \to X''$ ;

alors
\[
  \mathscr{G} \simeq \underline{\mathrm{Hom}}(\Omega^1_{X_0/T_0}, f_0^*(I))
  = \mathscr{V}_{X_0/T_0} \otimes f_0^*(\underline{I})
\]
$\mathscr{P}$ \struck{un} \add{$\mathscr{G}$-torseur à droite,}
[faisceau principal homogène \add{à droite} sous $\mathscr{G}$.]
\[
  \mathscr{P} \times \Omega^\bullet_{X'/T} \to \Omega^\bullet_{X''/T}
  \qquad \text{i.e.}
\]
\begin{align*}
  \mathscr{P} &\xrightarrow{u} \underline{\mathrm{Isom}}_{\text{de complexes}}
    (\Omega^\bullet_{X'/T}, \Omega^\bullet_{X''/T})
    \subset Z^0\, \underline{\mathrm{Hom}}^\bullet
    (\Omega^\bullet_{X'/T}, \Omega^\bullet_{X''/T}) \\
  \mathscr{G} &\xrightarrow{v} \underline{\mathrm{Aut}}_{\text{de complexes}}
    (\Omega^\bullet_{X''/T})
    \subset Z^0\, \underline{\mathrm{Hom}}^\bullet
    (\Omega^\bullet_{X''/T}, \Omega^\bullet_{X''/T}) \\
  \mathscr{G} &\xrightarrow{k} \underline{\mathrm{Hom}}^{(-1)}
    (\Omega^\bullet_{X''/T}, \Omega^\bullet_{X''/T})
\end{align*}
par produits intérieurs. On a
\[
  \begin{cases}
    u(pg) = u(p)\, v(g) \\
    v(g) - \mathrm{id} = d\, k(g) + k(g)\, d
  \end{cases}
  \qquad \text{si } p \in \Gamma(U, \mathscr{P}),\ g \in \Gamma(U, \mathscr{G})
\]
\note{« Isom » et « Aut » sont écrits sur un mot raturé ; la lettre lue $\mathscr{V}$ comme p. 229}

\page{234}
\note{p. « 2 » de l'auteur}
Plus généralement sur un espace \add{annelé} $X$, soient $K'$, $K''$ deux
complexes de Modules, et

$\mathscr{G}$ un faisceau abélien

$\mathscr{P}$ un $\mathscr{G}$-torseur à droite
\begin{align*}
  \mathscr{P} &\xrightarrow{u} \underline{\mathrm{Hom}}^0(K', K'') \\
  \mathscr{G} &\xrightarrow{v} \underline{\mathrm{Hom}}^\bullet(K'', K'') \\
  \mathscr{G} &\xrightarrow{k} \underline{\mathrm{Hom}}^{-1}(K', K'') ,
\end{align*}
\note{les symboles « Hom » de la première et de la troisième ligne sont écrits sur une lettre raturée ; le premier argument de la troisième ligne est lu $K'$}
\begin{gather*}
  u(pg) = u(p)\, v(g) \\
  v(g + g') = v(g) + v(g') \\
  v(g) = d\, k(g) + k(g)\, d
\end{gather*}
On veut un \add{canonique} homom.\ dans la cat.\ dérivée
\[
  K' \to K''
\]
d'où
\[
  \mathbb{R}\Gamma_X(K') \to \mathbb{R}\Gamma_X(K'')
\]
(compatible avec la structure de \struck{\ill{}} $\Gamma_X(\mathcal{O}_X)$ …)
[et plus gén.\ en relativisant par rapport à un morphisme d'espaces annelés
$X \to T$ …]

On introduit l'extension $\overline{\mathscr{G}}$ de $\mathbb{Z}$ par
$\mathscr{G}$
\[
  0 \to \mathscr{G} \xrightarrow{i} \overline{\mathscr{G}}
  \xrightarrow{j} \mathbb{Z} \to 0
\]
définie en termes de $\mathscr{P}$ par la relation
\[
  j^{-1}(1) \simeq \mathscr{P}
\]
(isom.\ de \struck{$\mathscr{G}$-is} $\mathscr{G}$-torseurs …).

Soit
\[
  \mathscr{H} = \underline{\mathrm{Hom}}^\bullet(K', K'')
\]
et soit $\mathscr{L}$ le complexe
\[
  \cdots \to 0 \to \mathscr{G} \to \overline{\mathscr{G}} \to 0 \to 0
\]
\note{au-dessus de $\mathscr{G}$ et de $\overline{\mathscr{G}}$, les indications de degré, entre guillemets, sont mal lisibles}

\page{235}
On trouve donc un homomorphisme
\[
  \varphi : \mathscr{L}^\bullet \to \mathscr{H}^\bullet
\]
par

\begin{tikzcd}[column sep=small]
  \cdots \arrow[r] & 0 \arrow[r] & \mathscr{G} \arrow[r, "i"] \arrow[d, "\varphi^{-1}"]
    & \overline{\mathscr{G}} \arrow[r] \arrow[d, "\varphi^0"] & 0 \arrow[r] & \cdots \\
  & \mathscr{H}^{-2} \arrow[r] & \mathscr{H}^{-1} \arrow[r, "\delta"]
    & \mathscr{H}^0 \arrow[r] & \mathscr{H}^1 \arrow[r] & \cdots
\end{tikzcd}

défini par les conditions
\[
  \varphi^{-1} = k , \qquad \varphi^0 | \mathscr{P} = u
\]
et la commutativité du carré.

[\struck{Def. \ill{} par} Un homom.\ de $\overline{\mathscr{G}}$ dans un
\uncertain{un} faisceau $\mathscr{F}$ = $\{$homom.\ $v$ de $\mathscr{G}$ dans
$\mathscr{F}$, et homom.\ $u$ de $\mathscr{P}$ dans $\mathscr{F}$, avec
condition $u(pg) = u(p)\, v(g)\}$. Donc prenant $\mathscr{F} = \mathscr{H}^0$,
et $u$, $v$ comme dessus, on trouve $\varphi^0$, avec $\varphi^0 i = u$. Mais
$u = \delta$\struck{$\varphi$}$k = \delta \varphi^{-1}$ …]
\note{ici $\mathscr{F}$ désigne un faisceau quelconque, non celui des pp. 223 à 227 ; les deux « $u$ » de la dernière ligne sont écrits tels quels}

Tensorisant avec $\mathcal{O}_X$, on trouve
\[
  \underline{\mathcal{O}}_X \xrightarrow{\ \sim\ } \text{\struck{$\mathscr{L}$}}
  \ \mathscr{L} \to \mathscr{H} \to \mathbb{R}\,\underline{\mathrm{Hom}}(K', K'')
\]
\add{dans cat.\ dérivée} \struck{On trouve donc} [\add{donc}
$\underline{\mathcal{O}}_X \to \mathbb{R}\,\underline{\mathrm{Hom}}(K', K'')$ dans
cat.\ dérivée]
\note{la flèche $\simeq$ porte « dans cat. dérivée » ; le terme intermédiaire est écrit sur une rature}

Donc
\[
  \text{\struck{\ill{}}}\ \mathbb{R}\Gamma_X(\underline{\mathcal{O}}_X)
  \to \text{\struck{\ill{}}}\ \mathbb{R}\,\underline{\mathrm{Hom}}(K', K'') ,
\]
d'où
\[
  \boxed{\; \mathbb{R}\Gamma_X(\underline{\mathcal{O}}_X) \to
  \mathbb{R}\,\underline{\mathrm{Hom}}(K', K'') \;}
\]
Par suite, \struck{faisant} passant au $\mathbb{R}\Gamma^0_X$, on trouve
\add{un élément de} $\mathbb{R}^0\underline{\mathrm{Hom}}(K', K'')$, i.e.\ un
homom.\ de la \add{cat.\ dérivée}
\note{la phrase se poursuit au-delà de la fin du lot}

\end{document}
